% Desarrollo reunido de DESARROLLO_MATEMATICO, apartados 38--40 y 50.
\subsubsection{Indice de capacité, lacunes et induction du régime tritique}
\label{vac-capacidad-trit}

La suite des lacunes est liée au raffinement par une
identité d'incréments. Pour la montrer, on fixe l'origine de phase
\(\theta=0\). Soient
\[
 \rho=\log_{10}9=1-\delta,\qquad
 \mathcal C_t=\lfloor(t+1)\rho\rfloor\quad(t\geq0).
\]
L'indice \(\mathcal C_t\) exprime la capacité décimale du raffinement
dans la normalisation qui compare un bloc ternaire de \(6\) positions,
avec \(729\) possibilités, et un bloc décimal de \(3\) positions, avec
\(1000\) possibilités. En effet,
\(\log_{1000}729=\log_{10}9=\rho\). La lettre \(K\) est réservée au
registre dodécaphasique qui sera construit ultérieurement.

\begin{proposition}[Complémentarité de la capacité et de la lacune]
\label{prop:capacidad-vacancia}
Pour \(t\geq1\),
\begin{equation}
 \mathcal C_t-\mathcal C_{t-1}=1-z_t(0).
 \label{eq:capacidad-vacancia}
\end{equation}
Par conséquent, si \(0\leq a<b\),
\[
 \mathcal C_b-\mathcal C_a+
 \sum_{t=a+1}^{b}z_t(0)=b-a.
\]
\end{proposition}
\begin{proof}
L'irrationalité de \(\delta\) donne
\[
 \mathcal C_t
 =t+1-\lceil(t+1)\delta\rceil
 =t-\lfloor(t+1)\delta\rfloor.
\]
La soustraction de deux termes consécutifs produit
\eqref{eq:capacidad-vacancia}. La somme sur l'intervalle est télescopique.
\end{proof}

Chaque transition augmente d'une unité le compteur des blocs. Si
\(z_t=0\), la capacité augmente également ; si \(z_t=1\), la capacité se
maintient tandis qu'une transition est incorporée à l'histoire. La lacune
est ici un événement de rétention de capacité, de position et
d'antécédent déterminés. Le bilan précédent conserve exactement la
longueur entre capacité acquise et lacunes enregistrées. La phase
observée \(g_t^{\rm cap}=[\mathcal C_t]_9\) vérifie
\[
 g_t^{\rm cap}-g_{t-1}^{\rm cap}\equiv1-z_t\pmod9.
\]
Cette phase de capacité et la phase chronologique \([t]_9\) sont deux
observations différentes. Une rétention locale de capacité et un tour
complet de l'holonomie ont des mécanismes distincts, bien que les deux
permettent de conserver l'observation de phase avec modification de l'état.

La relation avec TRIT peut se suivre dans les positions mêmes des lacunes.
Pour \(j\geq1\), définissons
\[
 p_j=\left\lfloor\frac j\delta\right\rfloor,\qquad
 v_j=\mathcal C_{p_j},\qquad h_j=p_{j+1}-p_j.
\]
L'inégalité \(p_j\delta<j<(p_j+1)\delta\) identifie \(p_j\)
comme la position de l'événement \(j\). Comme \(0<\delta<1\), elle donne aussi
\(\lfloor(p_j+1)\delta\rfloor=j\), et par conséquent
\[
 v_j=p_j-j,\qquad
 v_{j+1}-v_j=h_j-1.
\]
Si \(s_j=22-h_j\), alors \(s_j\in\{0,1\}\) et
\begin{equation}
 [v_{j+1}]_3=[v_j-s_j]_3.
 \label{eq:vacancia-regimen}
\end{equation}
Un intervalle de longueur \(22\) conserve la classe ; un intervalle de longueur
\(21\) la décale d'une unité dans le sens négatif. Le sélecteur canonique
s'applique à la phase positive \(r_j=1+[v_j]_9\) :
\[
 \tau_j=M_{\rm ph}(r_j).
\]
Un régime tritique est ainsi déterminé par l'indice de capacité
retenu à chaque événement, en plus du codage binaire de la lacune.

Les positions des intervalles courts forment le codage mécanique
induit de pente \(22-1/\delta\). Leurs espacements \(6\) et \(7\)
déterminent les longueurs des plateaux de la classe \([v_j]_3\), une
fois l'origine et la convention des extrémités fixées. Les premières classes
sont \(2\) pendant \(6\) événements, \(1\) pendant \(7\) et \(0\) pendant
\(7\) ; le premier segment de \(6\) est la restriction initiale d'un plateau.
Leurs signatures parcourent \(0,-1,+1\). Les équations
\eqref{eq:capacidad-vacancia} et \eqref{eq:vacancia-regimen} explicitent le
lien entre raffinement, lacunes et sélection du régime. La lecture
quadratique de TRIT réalise ensuite ces types au moyen de leurs générateurs
respectifs ; l'événement temporel et sa représentation géométrique
conservent ainsi une dépendance identifiable.

L'inversion de l'indice monotone fournit une autre observation utile :
\[
 \mathcal T(k)=\min\{t:\mathcal C_t=k\}
   =\left\lceil\frac k\rho\right\rceil-1\qquad(k\geq1).
\]
Si \(\sigma=\rho^{-1}-1\), le nombre de blocs nécessaires pour acquérir
\(m\) unités de capacité est
\[
 \mathcal T(k+m)-\mathcal T(k)
 =m+\lceil(k+m)\sigma\rceil-\lceil k\sigma\rceil
 \in\{m+\lfloor m\sigma\rfloor,m+\lceil m\sigma\rceil\}.
\]
En particulier, \(9\) unités de capacité nécessitent \(9\) ou \(10\)
blocs. La période d'une phase finie reste exacte, tandis que le
temps de retour mesuré par l'autre indice admet deux valeurs. Cette
distinction fait partie de la dynamique du raffinement et ne constitue pas une correction
de ses valeurs numériques a posteriori.
